\documentclass{article}
\usepackage[a4paper,margin=2.4cm]{geometry}
\usepackage{graphicx}
\usepackage{float}
\usepackage{fancyvrb}
\fvset{frame=single,fontsize=\small,framesep=3mm}

\title{Labwork 7: Reduction}
\author{Pham Tien Nhat}
\date{\today}

\begin{document}

\maketitle

\section{Implementation}

I implemented the grayscale stretch with 3 steps: convert the image to gray, find the max and min intensity, and recalculate every pixel from $[min, max]$ to $[0, 255]$.

The input image already uses the  range $[0, 255]$, so it is squeezed $[60, 160]$. 

\begin{Verbatim}
import math
import time
import numpy as np
import matplotlib.pyplot as plt

image = plt.imread('1.jpeg')
image = (image * (100 / 255) + 60).astype(np.uint8)
plt.imshow(image)
h, w, d = image.shape
\end{Verbatim}
\begin{figure}[H]
\centering
\includegraphics[width=0.6\textwidth]{../../image/lab7_input.png}
\caption{Input image, squeezed to $[60, 160]$.}
\label{figin}
\end{figure}

\textbf{Grayscale:} Each thread converts one pixel to gray with the same 2D kernel taken from previous labworks.

\begin{Verbatim}
@cuda.jit
def grayscale(src, dst):
    tidx = cuda.threadIdx.x + cuda.blockIdx.x * cuda.blockDim.x
    tidy = cuda.threadIdx.y + cuda.blockIdx.y * cuda.blockDim.y
    if tidx >= src.shape[1] or tidy >= src.shape[0]:
        return
    dst[tidy, tidx] = np.uint8((src[tidy, tidx, 0] + src[tidy, tidx, 1]
                                + src[tidy, tidx, 2]) / 3)

blockSize = (32, 32)
gridSize = (math.ceil(w / blockSize[0]), math.ceil(h / blockSize[1]))
devSrc = cuda.to_device(image)
devGray = cuda.device_array((h, w), np.uint8)
grayscale[gridSize, blockSize](devSrc, devGray)
gray = devGray.copy_to_host()
\end{Verbatim}
\begin{figure}[H]
\centering
\includegraphics[width=0.6\textwidth]{../../image/lab7_gray.png}
\caption{Gray image, input of step 2 and step 3.}
\label{figgray}
\end{figure}

\textbf{Find max/min:} The kernel follows the \texttt{sum1D} kernel from the slide. 

\begin{Verbatim}
sharedSize = 256

@cuda.jit
def reduceMax(src, dst):
    cache = cuda.shared.array((sharedSize, ), np.uint8)
    localtid = cuda.threadIdx.x
    tid = cuda.threadIdx.x + cuda.blockIdx.x * cuda.blockDim.x
    if tid < src.shape[0]:
        cache[localtid] = src[tid]
    else:
        cache[localtid] = 0
    cuda.syncthreads()
    s = 1
    while s < cuda.blockDim.x:
        if localtid % (s * 2) == 0:
            cache[localtid] = max(cache[localtid], cache[localtid + s])
        cuda.syncthreads()
        s = s * 2
    if localtid == 0:
        dst[cuda.blockIdx.x] = cache[0]
\end{Verbatim}

\texttt{reduceMin} is the same kernel with \texttt{min} and a padding value of 255.


\begin{Verbatim}
devData = devGray.reshape(h * w)
gridReduce = math.ceil(h * w / sharedSize)
devMax = cuda.device_array(gridReduce, np.uint8)
devMin = cuda.device_array(gridReduce, np.uint8)
reduceMax[gridReduce, sharedSize](devData, devMax)
reduceMin[gridReduce, sharedSize](devData, devMin)

gmax = devMax.copy_to_host().max()
gmin = devMin.copy_to_host().min()
\end{Verbatim}

I get the values: $max = 160$ and $min = 60$.

\textbf{Stretch: } Each pixel is recalculated with $g' = \frac{g - min}{max - min} \times 255$. 

\begin{Verbatim}
@cuda.jit
def stretch(src, dst, gmin, gmax):
    tidx = cuda.threadIdx.x + cuda.blockIdx.x * cuda.blockDim.x
    tidy = cuda.threadIdx.y + cuda.blockIdx.y * cuda.blockDim.y
    if tidx >= src.shape[1] or tidy >= src.shape[0]:
        return
    dst[tidy, tidx] = np.uint8((src[tidy, tidx] - gmin)
                               / (gmax - gmin) * 255)

devStretch = cuda.device_array((h, w), np.uint8)
stretch[gridSize, blockSize](devGray, devStretch,
                             np.float32(gmin), np.float32(gmax))
stretched = devStretch.copy_to_host()
\end{Verbatim}
\begin{figure}[H]
\centering
\includegraphics[width=0.9\textwidth]{../../image/lab7_stretch.png}
\caption{Gray image $[60, 160]$ (left) and stretched image $[0, 255]$ (right).}
\label{figstretch}
\end{figure}

Figure~\ref{fighist} shows the histogram: the gray values are packed in $[60, 160]$ before the stretch, and scatter over $[0, 255]$ , showing cleaner contrast.

\begin{figure}[H]
\centering
\includegraphics[width=0.65\textwidth]{../../image/lab7_hist.png}
\caption{Histogram before and after the stretch.}
\label{fighist}
\end{figure}

\section{Experimental Results}
4 2D block sizes are tested for the two map kernels: $(4, 4)$, $(8, 8)$, $(16, 16)$ and $(32, 32)$.

\begin{Verbatim}
block_sizes = [(4, 4), (8, 8), (16, 16), (32, 32)]
times_gray = []
times_stretch = []

for blockSize in block_sizes:
    gridSize = (math.ceil(w / blockSize[0]), math.ceil(h / blockSize[1]))

    start = time.time()
    grayscale[gridSize, blockSize](devSrc, devGray)
    cuda.synchronize()
    times_gray.append(time.time() - start)

    start = time.time()
    stretch[gridSize, blockSize](devGray, devStretch,
                                 np.float32(gmin), np.float32(gmax))
    cuda.synchronize()
    times_stretch.append(time.time() - start)

plt.figure(figsize=(8, 5))
plt.plot([str(b) for b in block_sizes], times_gray, marker='o',
         label='grayscale')
plt.plot([str(b) for b in block_sizes], times_stretch, marker='o',
         label='stretch')
plt.xlabel("2D Block Size")
plt.ylabel("Time (seconds)")
plt.title("CUDA 2D Block Size vs GPU Execution Time")
plt.legend()
plt.grid(True)
plt.show()
\end{Verbatim}

Figure~\ref{fig7} shows the relationship between the 2D block size and the GPU time.

\begin{figure}[H]
\centering
\includegraphics[width=0.7\textwidth]{../../image/lab7.png}
\caption{2D block size vs GPU time.}
\label{fig7}
\end{figure}

The block of $(4, 4)$ is the slowest for both kernels. $(32, 32)$ is the fastest where the peak of stretch comes at $(16, 16)$.

\end{document}
